\documentclass[12pt,a4paper]{article}

\usepackage[top=1in,bottom=1in,left=1.1in,right=1.1in]{geometry}
\usepackage[T1]{fontenc}
\usepackage[utf8]{inputenc}
\usepackage{lmodern}
\usepackage{microtype}
\usepackage{parskip}
\usepackage{array}
\usepackage{booktabs}
\usepackage{longtable}
\usepackage{tabularx}
\usepackage{enumitem}
\usepackage{graphicx}
\usepackage{xcolor}
\usepackage{hyperref}
\usepackage{titlesec}
\usepackage{fancyhdr}
\usepackage{lastpage}
\usepackage{listings}
\usepackage{tikz}

\definecolor{DocBlue}{RGB}{34,52,74}
\definecolor{DocGray}{RGB}{90,90,90}
\definecolor{RuleGray}{RGB}{180,180,180}
\definecolor{CodeBack}{RGB}{247,247,247}
\definecolor{CodeFrame}{RGB}{210,210,210}

\hypersetup{
  colorlinks=true,
  linkcolor=DocBlue,
  urlcolor=DocBlue,
  citecolor=DocBlue,
  pdftitle={FlexQL Design Document}
}

\titleformat{\section}
  {\Large\bfseries\color{DocBlue}}
  {\thesection}{0.8em}{}
  [\color{RuleGray}\titlerule]

\titleformat{\subsection}
  {\large\bfseries\color{DocBlue}}
  {\thesubsection}{0.8em}{}

\titleformat{\subsubsection}
  {\normalsize\bfseries\color{DocGray}}
  {\thesubsubsection}{0.8em}{}

\lstdefinestyle{docstyle}{
  language=bash,
  basicstyle=\ttfamily\small,
  backgroundcolor=\color{CodeBack},
  frame=single,
  rulecolor=\color{CodeFrame},
  breaklines=true,
  showstringspaces=false,
  columns=fullflexible
}
\lstset{style=docstyle}

\pagestyle{fancy}
\fancyhf{}
\fancyhead[L]{\small\color{DocGray}FlexQL Design Document}
\fancyhead[R]{\small\color{DocGray}\leftmark}
\fancyfoot[C]{\small\color{DocGray}\thepage\ /\ \pageref{LastPage}}
\setlength{\headheight}{14pt}
\renewcommand{\headrulewidth}{0.3pt}
\renewcommand{\footrulewidth}{0pt}

\setlength{\parindent}{0pt}
\setlength{\parskip}{6pt}

\newcommand{\projecttitle}{FlexQL}
\newcommand{\projectsubtitle}{A Flexible SQL-like Database Driver}

\begin{document}

\begin{titlepage}
\thispagestyle{empty}

\begin{tikzpicture}[remember picture,overlay]
  \draw[line width=1pt,color=RuleGray]
    ([xshift=1cm,yshift=-1cm]current page.north west)
    rectangle
    ([xshift=-1cm,yshift=1cm]current page.south east);
  \draw[line width=2pt,color=DocBlue]
    ([xshift=1.35cm,yshift=-1.35cm]current page.north west)
    -- ([xshift=-1.35cm,yshift=-1.35cm]current page.north east);
\end{tikzpicture}

\vspace*{2.2cm}

\begin{center}
{\fontsize{28}{32}\selectfont\bfseries\color{DocBlue}\projecttitle\par}
\vspace{0.4cm}
{\Large\projectsubtitle\par}
\vspace{0.35cm}

\end{center}

\vspace{2cm}

\renewcommand{\arraystretch}{1.8}
\begin{center}
\begin{tabular}{>{\bfseries}p{4.2cm}p{8.5cm}}
Project Name & FlexQL \\
Name & Swagnik Ghosh \\
Roll No. & 25CS60R67 \\
Department & Computer Science and Engineering \\
Date & \today \\
\end{tabular}
\end{center}

\vfill

\begin{center}
\color{DocGray}
\small
This document presents the design, implementation decisions, and technical
structure of the FlexQL database project based on the current repository state.
\end{center}

\end{titlepage}

\tableofcontents
\newpage

\section{Project Overview}

FlexQL is a SQL-like relational database system implemented in C++17. The
project follows a client-server architecture in which a TCP server accepts SQL
queries, parses them, executes them against persistent storage, and returns
formatted results to connected clients.

The repository shows that the project includes:

\begin{itemize}[leftmargin=1.8em]
  \item a database server started using \texttt{./server},
  \item an interactive client / REPL (\texttt{bin/flexql-client}),
  \item benchmarking and unit-test execution started using \texttt{./benchmark}.
\end{itemize}

The design emphasizes simplicity, inspectable storage files, practical
performance for assignment-scale workloads, and modular separation between
parsing, execution, storage, indexing, concurrency, and networking.

\section{Objectives}

\begin{itemize}[leftmargin=1.8em]
  \item Implement a working SQL-like database from scratch in C++17.
  \item Support persistent storage for tables, schemas, and indexes.
  \item Provide efficient primary-key lookup and acceptable full-scan
        performance.
  \item Expose the system through a networked client-server interface.
  \item Support repeated testing and benchmarking of the complete system.
\end{itemize}

\section{Supported Features}

Based on the current codebase and README, FlexQL supports the following major
features:

\begin{itemize}[leftmargin=1.8em]
  \item \texttt{CREATE DATABASE} and \texttt{USE}
  \item \texttt{CREATE TABLE} and \texttt{DROP TABLE}
  \item \texttt{SHOW DATABASES} and \texttt{SHOW TABLES}
  \item \texttt{INSERT INTO ... VALUES (...)}
  \item optional \texttt{EXPIRES AT 'YYYY-MM-DD HH:MM:SS'} on insert
  \item \texttt{SELECT *} and projected-column selection
  \item \texttt{WHERE} filtering using equality and comparison operators
  \item single-table queries and \texttt{INNER JOIN}
  \item primary-key uniqueness enforcement
  \item query-result caching
  \item row expiration with lazy filtering
  \item benchmark and validation support
\end{itemize}

\subsection{How the Features Are Implemented}

\subsubsection{Database and Table Management}

Database creation is handled by creating the required on-disk directory layout
under \texttt{data/databases/<db>/}. The executor checks whether the database
already exists and supports \texttt{IF NOT EXISTS} behavior for safe repeated
execution.

Table creation parses the declared schema, validates supported data types, and
persists the metadata to a separate \texttt{.schema} file. A corresponding data
file and index state are then initialized so the table is immediately ready for
insertion and lookup.

Dropping a table removes the table data file, schema file, index snapshot,
delta log, and bloom filter file. This keeps logical deletion simple and avoids
stale table metadata remaining on disk.

\subsubsection{SQL Parsing}

The parser in \texttt{src/parser/sql\_parser.cpp} uses statement-type dispatch
based on the SQL prefix and then applies targeted parsing logic for each
statement family. It supports:

\begin{itemize}[leftmargin=1.8em]
  \item qualified table names such as \texttt{db.table},
  \item inline and table-level \texttt{PRIMARY KEY} syntax,
  \item multi-row \texttt{INSERT} statements,
  \item \texttt{EXPIRES AT} for single-row inserts,
  \item comparison operators in \texttt{WHERE} clauses,
  \item join specifications for \texttt{INNER JOIN}.
\end{itemize}

The parser normalizes identifiers and builds a structured \texttt{Query}
object, which is then passed to the executor.

\subsubsection{Insert Processing}

Insert execution follows a clear pipeline:

\begin{enumerate}[leftmargin=1.8em]
  \item parse the statement into a \texttt{Query} object,
  \item resolve the active database,
  \item validate the values against the schema,
  \item compute the expiration time using either \texttt{EXPIRES AT} or the
        default TTL,
  \item serialize the row into the append-only table file,
  \item update the in-memory primary-key index,
  \item invalidate dependent cache entries.
\end{enumerate}

The code supports both single-row and multi-row insertion. For multi-row
inserts, rows are prepared in memory first, duplicate primary keys are checked
within the batch and against existing data, and then the combined payload is
written efficiently in one append path.

\subsubsection{Select Queries}

Simple \texttt{SELECT} queries are executed using two main paths:

\begin{itemize}[leftmargin=1.8em]
  \item if the \texttt{WHERE} clause is an equality condition on the primary
        key, the executor performs a direct index lookup and reads the row by
        file offset;
  \item otherwise, the executor performs a scan-based path over cached row
        pages loaded from the table.
\end{itemize}

This is an important optimization because it gives fast point lookup without
making the implementation overly complex.

\subsubsection{Filtering and Comparison}

The parser supports \texttt{=}, \texttt{!=}, \texttt{<>}, \texttt{>},
\texttt{>=}, \texttt{<}, and \texttt{<=} in \texttt{WHERE} conditions. During
execution, the relevant column is resolved from the schema and values are
compared according to the declared column type. This allows the same filtering
logic to work across integer, decimal, string, and datetime-like values.

\subsubsection{Join Execution}

Join queries are implemented in the executor using a hash-join style approach.
The left table is read first, then the right-side rows are loaded and hashed on
the join key. Matching rows are combined into joined records, and projection or
post-join filtering is applied before results are returned.

The implementation also resolves qualified and unqualified column references,
detects ambiguous columns, and supports chained joins across multiple tables.

\subsubsection{Caching}

FlexQL uses caching at two levels:

\begin{itemize}[leftmargin=1.8em]
  \item a query-result cache for repeated \texttt{SELECT} statements,
  \item a page cache that stores 64-row table pages after an initial scan.
\end{itemize}

The query cache uses a segmented LRU structure with probation and protected
regions. This means frequently reused queries survive longer than one-time
queries. Cache invalidation is triggered after writes so stale results are not
served.

\subsubsection{Expiration and TTL}

Each stored row carries an expiration timestamp. If the insert statement does
not specify \texttt{EXPIRES AT}, the executor uses the \texttt{TtlManager} to
assign a default expiration window. On reads, expired rows are skipped rather
than physically deleted.

This implementation keeps insertion simple while still giving the system a
useful row-lifetime feature.

\subsubsection{Networking and Client Interaction}

Clients communicate with the server using framed TCP messages. The server
parses a request, executes it through the shared executor logic, and returns
either an \texttt{OK}, \texttt{ERR}, or tabular \texttt{RESULT} response.

The REPL builds on the same client API and adds:

\begin{itemize}[leftmargin=1.8em]
  \item pretty-printed table output,
  \item per-query execution timing.
\end{itemize}

\subsubsection{Testing and Benchmarking Support}

The project uses a two-terminal workflow in which the server is started first
and unit tests or benchmarks are then started from another terminal through the
\texttt{./benchmark} wrapper. This is an important implemented feature because
it shows that the system was built not only to run queries, but also to be
validated and measured in a repeatable way.

\section{High-Level Architecture}

The project is organized into clear modules under \texttt{include/} and
\texttt{src/}. The implementation follows the layered flow shown below.

\begin{center}
\renewcommand{\arraystretch}{1.35}
\begin{tabular}{>{\bfseries}p{4.4cm}p{8.4cm}}
\toprule
Layer & Responsibility \\
\midrule
Client Layer & REPL client, benchmark client, public API wrappers \\
Network Layer & TCP communication, message framing, server back-ends \\
Parser Layer & SQL parsing into internal query structures \\
Execution Layer & Validation, planning decisions, scans, joins, caching \\
Storage Layer & Table files, schema files, row serialization, mmap reads \\
Index Layer & Primary-key index, bloom filter, snapshot and delta persistence \\
Concurrency Layer & Table locking, append coordination, thread pool support \\
\bottomrule
\end{tabular}
\end{center}

This organization keeps major responsibilities separate and makes the project
easier to reason about, test, and extend.

\section{Repository Structure}

\begin{lstlisting}[caption={Relevant repository structure}]
include/
src/
  client/
  server/
  parser/
  query/
  storage/
  index/
  cache/
  concurrency/
  expiration/
  network/
tests/
bin/
data/databases/
\end{lstlisting}

\subsection{Important Source Areas}

\begin{itemize}[leftmargin=1.8em]
  \item \texttt{src/server/}: server entry point and back-end selection
  \item \texttt{src/client/}: client API, REPL, and benchmark client
  \item \texttt{src/parser/}: SQL parser implementation
  \item \texttt{src/query/}: executor, join handling, vectorized comparison
  \item \texttt{src/storage/}: table, schema, row encoding, mmap reader
  \item \texttt{src/index/}: persisted primary-key index and bloom filter
  \item \texttt{src/concurrency/}: lock manager, lock-free append, thread pool
\end{itemize}

\section{Data Model and Type System}

FlexQL stores rows in a row-oriented form. Each table has a schema file
describing the columns and one primary key used for indexed lookup.

\subsection{Supported Types}

\begin{center}
\renewcommand{\arraystretch}{1.3}
\begin{tabular}{lll}
\toprule
Type & SQL Form & Notes \\
\midrule
Integer & \texttt{INT} & validated on insert \\
Decimal & \texttt{DECIMAL(p,s)} & stored as text, parsed when needed \\
String & \texttt{VARCHAR(n)} & stored as text \\
Date-Time & \texttt{DATETIME} & used for timestamp-like values \\
\bottomrule
\end{tabular}
\end{center}

The implementation stores row values as strings internally for simplicity.
Validation happens at insertion time based on the declared schema.

\subsection{Core Runtime Structures}

The codebase uses data structures such as \texttt{Row}, \texttt{QueryResult},
and parsed \texttt{Query} objects to carry data across modules. This keeps the
pipeline clear: parse SQL, build a query object, execute it, and return a
structured result to the client.

\subsubsection{Query Structure}

The \texttt{Query} structure is the central intermediate representation used by
the parser and executor. It stores:

\begin{itemize}[leftmargin=1.8em]
  \item the statement type,
  \item database and table names,
  \item \texttt{IF EXISTS} and \texttt{IF NOT EXISTS} flags,
  \item column definitions for table creation,
  \item single-row or multi-row insert payloads,
  \item expiration metadata,
  \item projected columns for \texttt{SELECT},
  \item a parsed \texttt{WHERE} condition,
  \item join specifications.
\end{itemize}

This design avoids repeatedly reparsing SQL during execution because all
necessary information is collected into one object before the executor begins
its work.

\subsubsection{Result Representation}

\texttt{QueryResult} is the unified response object used across the server. It
contains a success flag, a message string, result column names, result rows,
and a cache expiration timestamp. This is useful because both DDL/DML commands
and table-returning queries share the same response type, simplifying response
building in the server and client layers.

\section{Storage Design}

\subsection{Physical Layout}

Each database is stored under:

\begin{center}
\texttt{data/databases/<db>/}
\end{center}

Within each database, FlexQL stores tables and indexes separately:

\begin{center}
\renewcommand{\arraystretch}{1.3}
\begin{tabular}{lll}
\toprule
File & Purpose & Example \\
\midrule
\texttt{tables/<table>.tbl} & row data & \texttt{student.tbl} \\
\texttt{tables/<table>.schema} & schema metadata & \texttt{student.schema} \\
\texttt{indexes/<table>.idx} & index snapshot & \texttt{student.idx} \\
\texttt{indexes/<table>.idx.delta} & pending index updates & \texttt{student.idx.delta} \\
\texttt{indexes/<table>.idx.bloom} & bloom filter data & \texttt{student.idx.bloom} \\
\bottomrule
\end{tabular}
\end{center}

\subsection{Row Format}

Rows are serialized into newline-terminated text records with an expiration
field followed by column values. Values are escaped before persistence so that
the delimiter format remains parseable.

\begin{center}
\texttt{expiration|value1|value2|...|valueN}
\end{center}

This design keeps table files easy to inspect and debug while still supporting
persistent storage.

\subsection{Append-Only Writes}

Table insertion is append-only. The implementation in
\texttt{src/concurrency/lockfree\_append.cpp} reserves file offsets
atomically and uses positional writes. This avoids rewriting existing rows and
helps concurrent inserts remain simple and safe.

In the \texttt{Table} class, insertion is split into multiple steps:

\begin{enumerate}[leftmargin=1.8em]
  \item ensure the table and index are loaded,
  \item validate incoming values against the schema,
  \item normalize quoted literals,
  \item serialize the row into the text row format,
  \item reserve byte space in the append file,
  \item write the bytes using positional write operations,
  \item update the in-memory primary-key index,
  \item flush index delta entries when the pending threshold is reached.
\end{enumerate}

To prevent duplicate primary keys under concurrent insert activity, the table
tracks both committed index entries and an \texttt{inflight\_primary\_keys\_}
set. This is an important implementation detail because it prevents two writers
from reserving space for the same primary key before the index update becomes
visible.

\subsubsection{Single-Row and Batch Inserts}

The project implements both \texttt{insert\_row()} and \texttt{insert\_rows()}.
The single-row path is straightforward and optimized for regular SQL usage.
The multi-row path prepares all rows first, validates duplicate keys inside the
batch, computes offsets for the full batch, concatenates the serialized
records, and writes them in one larger append. This reduces syscall overhead
and improves throughput during benchmark workloads.

\subsection{Memory-Mapped Reads}

The read path uses \texttt{mmap} through \texttt{MmapReader}. This allows the
system to scan table files efficiently and reduces unnecessary copying during
full-table reads.

The \texttt{Table} class uses a fast path and a slow path for reads:

\begin{itemize}[leftmargin=1.8em]
  \item if the reader is already loaded and the mapping is not stale, read
        operations proceed immediately under a shared lock;
  \item if the file has grown and the mapped view is stale, the table refreshes
        the mapping under a unique lock and then resumes reading.
\end{itemize}

This keeps common reads lightweight while still handling append growth safely.

\subsubsection{Materialized Rows and Row Views}

The storage code supports both fully materialized \texttt{Row} objects and
zero-copy \texttt{RowView} records. \texttt{RowView} stores
\texttt{string\_view} slices into the memory-mapped file, which allows scan
paths to avoid allocating strings for rows that are filtered out. This is one
of the more performance-aware parts of the implementation.

\section{Schema Management}

Schema metadata is handled by the \texttt{Schema} module. It stores column
definitions, type information, and primary-key configuration.

Primary-key behavior in the current project is important because:

\begin{itemize}[leftmargin=1.8em]
  \item each table is expected to have a usable primary key,
  \item duplicate primary keys are rejected during insert,
  \item the primary key drives direct index-based lookup.
\end{itemize}

The schema is persisted in a separate human-readable \texttt{.schema} file,
which makes table metadata easy to inspect.

\subsection{Schema Validation Logic}

When a table is created, the parser first converts SQL type declarations into
internal \texttt{DataType} values. The schema layer then stores the ordered
column list and records which column is the primary key.

During insert, every row is validated column-by-column:

\begin{itemize}[leftmargin=1.8em]
  \item row width must match schema width,
  \item each value must satisfy the declared type,
  \item the primary-key column must be valid and unique.
\end{itemize}

This keeps the validation policy centralized in the storage path rather than
scattering it across multiple modules.

\section{Indexing Design}

\subsection{Primary-Key Index}

Despite the class name \texttt{BTreeIndex}, the current implementation uses an
in-memory hash map for primary-key to row-offset mapping. This is appropriate
for the workload pattern targeted by the project because equality lookup on the
primary key is the main fast path.

Internally, the index maps:

\begin{center}
\texttt{primary\_key -> table\_file\_byte\_offset}
\end{center}

This means a successful primary-key search does not scan the table. Instead,
the executor asks the table for the offset, then reads the row directly from
the mapped table file.

\subsection{Index Persistence}

Index data is persisted using two complementary mechanisms:

\begin{itemize}[leftmargin=1.8em]
  \item a full snapshot file (\texttt{.idx}),
  \item an incremental delta log (\texttt{.idx.delta}).
\end{itemize}

This allows inserts to remain lightweight while still preserving index state
across restarts. At shutdown or compaction time, the full snapshot can be
rewritten and the delta log can be cleared.

The implementation uses the following persistence strategy:

\begin{itemize}[leftmargin=1.8em]
  \item on normal insert activity, new index entries accumulate in memory,
  \item once the pending-update threshold is reached, a drained batch is
        appended to the delta log,
  \item on table destruction or compaction, a full snapshot is written and the
        delta log is truncated,
  \item on startup, the snapshot is loaded first and then delta entries are
        replayed.
\end{itemize}

This gives better insert behavior than rewriting the full index file for every
new row.

\subsection{Bloom Filter}

The index layer also includes a bloom filter to reduce unnecessary hash-map
lookups for keys that are definitely absent. This is a small but useful
optimization for negative primary-key probes.

The bloom filter is updated when new keys are inserted, rebuilt when the index
is compacted, and loaded again during startup if the persisted bloom file is
available.

\section{Query Processing}

\subsection{Parsing}

The parser converts SQL text into an internal \texttt{Query} representation.
It recognizes the project's supported statements and validates the general
structure before execution.

The implementation uses an efficient dispatch pattern. Instead of applying one
large regex to every statement, it first normalizes the SQL prefix and chooses
the correct parsing routine based on the opening keywords. This reduces parsing
overhead, especially for large insert statements.

\subsection{Execution Paths}

The executor handles different query paths depending on the statement type:

\begin{itemize}[leftmargin=1.8em]
  \item DDL operations such as creating databases or tables
  \item insert operations with validation and persistence
  \item primary-key point lookup using the index
  \item full-table scan when no direct indexed path is available
  \item join execution for supported \texttt{INNER JOIN} queries
\end{itemize}

The executor also performs cache invalidation and expiration checking so that
query results stay consistent with recent writes and TTL behavior.

\subsubsection{Executor Responsibilities}

The \texttt{Executor} class is the main orchestration component. It owns the
parser, cache reference, TTL manager, table cache, and page cache. Its
responsibilities include:

\begin{itemize}[leftmargin=1.8em]
  \item parsing raw SQL into a query object,
  \item resolving the active database,
  \item selecting the correct execution function,
  \item loading and caching table instances,
  \item updating query cache entries and invalidating them after writes,
  \item returning a unified \texttt{QueryResult}.
\end{itemize}

\subsubsection{Simple Select Path}

For simple selects, the executor first validates the projection and
\texttt{WHERE} column references. It then chooses between:

\begin{itemize}[leftmargin=1.8em]
  \item a primary-key equality fast path using the index and direct row read,
  \item a scan path using the page cache and row filtering.
\end{itemize}

For scan-based reads, the executor loads the table once into 64-row pages and
then reuses those pages for subsequent reads. This improves repeated-query
performance.

\subsection{Join Support}

FlexQL supports \texttt{INNER JOIN} and chained joins. The executor validates
qualified and unqualified column references and resolves ambiguity before
returning the final result set.

In the current implementation, joins are executed by:

\begin{enumerate}[leftmargin=1.8em]
  \item loading the left table rows,
  \item building a metadata list for joined columns,
  \item loading each right table,
  \item building an in-memory hash structure on the join key,
  \item probing that structure using rows from the partially joined result,
  \item constructing new joined rows for every match.
\end{enumerate}

This approach is practical for the project scope and supports chained joins by
repeating the same process across successive tables.

\subsection{Vectorized String Comparison}

The file \texttt{src/query/vectorized\_executor.cpp} includes an SSE-based
optimized equality path for x86\_64 systems. This is a targeted low-level
optimization used for string comparison when supported by the platform.

\section{Caching and Expiration}

\subsection{Query Cache}

The project includes a segmented LRU cache for query results. New entries are
placed in a probation segment and promoted to a protected segment on reuse.
This is more resistant to cache pollution than a single plain LRU list.

The implementation stores dependency information alongside cache entries. Each
cached query can be associated with one or more table dependency keys such as
\texttt{database.table}. When a table changes, the cache can invalidate only
the affected entries rather than clearing the entire cache every time.

\subsection{Page Cache}

The executor also maintains a row-page cache for table data, using 64-row page
grouping to reduce repeated reconstruction work during repeated scans.

This page cache is separate from the query-result cache. It stores decoded row
pages by page identifier, allowing future scans or joins to reuse loaded rows
without rereading and reparsing the table file from scratch.

\subsection{TTL and Expiration}

The TTL manager computes expiration timestamps and checks whether rows are
expired at read time. The current implementation applies a default TTL through
the \texttt{TtlManager}, while \texttt{EXPIRES AT} allows the caller to set an
explicit expiration timestamp for an inserted row.

Expired rows remain on disk and are filtered lazily during reads, which keeps
the write path simple.

Cached results also track the earliest relevant expiration time. This lets the
query cache reject entries whose logical freshness has passed even if the table
has not been modified.

\section{Concurrency Design}

FlexQL uses a practical concurrency model centered on table-level coordination.

\begin{itemize}[leftmargin=1.8em]
  \item A lock manager provides per-table synchronization.
  \item Readers can proceed concurrently using shared locking.
  \item Writers coordinate insert and metadata-changing paths safely.
  \item Append reservation and positional writes reduce lock contention during
        concurrent inserts.
\end{itemize}

This design is simpler than full transactional concurrency control and is
appropriate for the scope of the project.

\subsection{Locking Strategy}

The code uses multiple layers of synchronization for different purposes:

\begin{itemize}[leftmargin=1.8em]
  \item a per-table \texttt{shared\_mutex} for readers and writers,
  \item a write-state mutex for index updates and inflight key tracking,
  \item a separate mutex for serializing delta-log flushes.
\end{itemize}

This separation is useful because it avoids making all write-related work
compete on one oversized lock.

\subsection{Thread Pool Support}

In default server mode, accepted client connections are handed to a thread
pool. This allows the server to cap worker creation while still processing
multiple clients concurrently.

\section{Network and Server Design}

\subsection{Wire Model}

Clients communicate with the server over TCP. The server receives framed
messages, executes the SQL, and returns a structured textual response.

The message format is length-prefixed. A 32-bit network-order payload size is
sent before the message body. This avoids relying on delimiters in the raw TCP
stream and ensures the receiver knows exactly how many bytes belong to one
request or one response.

The response protocol uses three major forms:

\begin{itemize}[leftmargin=1.8em]
  \item \texttt{OK<TAB>message} for successful non-tabular commands,
  \item \texttt{ERR<TAB>message} for failures,
  \item \texttt{RESULT<TAB>colcount<TAB>rowcount} followed by headers and rows
        for select-style outputs.
\end{itemize}

\subsection{Server Back-Ends}

The main server binary supports three execution modes based on the startup
argument:

\begin{itemize}[leftmargin=1.8em]
  \item default threaded mode using a thread pool,
  \item \texttt{--epoll} mode,
  \item \texttt{--io-uring} mode.
\end{itemize}

This is a strong design feature of the project because it allows the same
database engine to be exercised under multiple network I/O models.

\subsubsection{Threaded Mode}

The default mode accepts connections and delegates client handling to the
thread pool. Each worker reads framed requests, calls the request processor,
and sends framed responses back to the client.

\subsubsection{Epoll Mode}

The epoll back-end keeps client state objects and multiplexes readiness events
through an event loop. This implementation is more scalable for large numbers
of mostly idle connections and demonstrates an event-driven server design.

\subsubsection{io\_uring Mode}

The codebase also includes an \texttt{io\_uring}-based back-end, which shows an
experimental modern Linux asynchronous I/O path. This makes the project more
interesting from a systems perspective even though the logical database engine
remains the same.

\section{Client API and User Interaction}

The project provides both programmatic and interactive access.

\begin{itemize}[leftmargin=1.8em]
  \item The client API exposes operations such as open, execute, and close.
  \item The REPL supports interactive SQL entry.
  \item The benchmark wrapper drives unit tests and performance measurement
        against a live server instance.
\end{itemize}

\subsection{C API Design}

The public C API provides an opaque \texttt{FlexQL} handle. The implementation
in \texttt{src/client/api.cpp} handles:

\begin{itemize}[leftmargin=1.8em]
  \item connection establishment,
  \item synchronized request submission,
  \item framed response reception,
  \item callback-based row delivery for result sets,
  \item error-string allocation for callers.
\end{itemize}

This design keeps the client library usable from C-style code while still
wrapping the C++ networking internals.

\subsection{REPL Behavior}

The REPL is implemented as a thin interactive layer above the client API. It
formats result sets into readable ASCII tables and reports execution time. This
makes the system practical both for manual testing and for demonstration.

\subsection{Batch Insert Protocol}

The client API also supports batch insertion through a special
\texttt{MSG\_BATCH\_INSERT} payload format. Multiple SQL insert statements are
packed into one framed request, and the server executes them sequentially as a
batch. This is especially useful for benchmark workloads.

\section{Build, Testing, and Benchmarking}

\subsection{Build}

\begin{lstlisting}[caption={Build command}]
make
\end{lstlisting}

The build generates the binaries used in the normal workflow.

\subsection{Testing}

The testing workflow used in this project is:

\begin{lstlisting}[caption={Unit-test workflow}]
make

# Terminal 1
./server

# Terminal 2
./benchmark --unit-test
\end{lstlisting}

This runs the unit-test mode after the server is started.

This is significant from an implementation perspective because it validates not
only the parser, but also the interaction between networking, execution,
storage, indexing, and expiration behavior.

\subsection{Benchmarking}

\begin{lstlisting}[caption={Benchmark workflow}]
make

# Terminal 1
./server

# Terminal 2
./benchmark
./benchmark 200000
\end{lstlisting}

The default command runs benchmark plus unit tests, and a custom row count can
be provided as a command-line argument.

This reflects the actual project workflow: the server runs first, and then the
benchmark is executed from another terminal against the live database service.

\section{Current Limitations}

From the repository and existing implementation, the main limitations are:

\begin{itemize}[leftmargin=1.8em]
  \item no full transactional engine or SQL standard compliance,
  \item no \texttt{UPDATE} or \texttt{DELETE} path in normal operation,
  \item row-oriented text storage is simpler but not space-optimal,
  \item append-only design causes stale or expired data to remain on disk,
  \item range-query support is limited because the fast path is hash-index
        based rather than tree-index based.
\end{itemize}

\section{Important Trade-Offs}

Instead of repeating trade-offs in every subsection, the most important design
trade-offs of FlexQL are summarized here.

\begin{center}
\renewcommand{\arraystretch}{1.35}
\begin{tabularx}{\textwidth}{>{\bfseries}p{3.2cm}X X}
\toprule
Decision & Benefit & Cost / Limitation \\
\midrule
Append-only storage &
simple writes, safer concurrent insert path, easy persistence model &
old or expired rows remain on disk and file size grows over time \\

Text-based row storage &
easy debugging and human-readable data files &
larger storage footprint and slower numeric processing than binary formats \\

Hash-based primary-key index &
very fast equality lookup on the main access path &
not ideal for ordered traversal or efficient range queries \\

Snapshot + delta index persistence &
cheap incremental index updates during insert-heavy workloads &
recovery logic is more complex than a single always-rewritten index file \\

Memory-mapped reads &
efficient scans with low-copy access patterns &
mapping lifecycle must be managed carefully when files grow \\

Segmented LRU caching &
better protection against one-time query pollution &
more implementation complexity than a plain LRU cache \\

Lazy TTL expiration &
keeps the write path simple and avoids physical deletion work &
expired rows still consume disk space until compaction or rebuild \\

Multiple server back-ends &
good experimental flexibility across threaded, epoll, and io\_uring modes &
increases implementation and maintenance complexity \\
\bottomrule
\end{tabularx}
\end{center}

\section{Conclusion}

FlexQL is a well-structured educational database project that combines core
database ideas with systems-level implementation in C++17. Its strongest
qualities are modular design, persistent storage, indexed primary-key lookup,
multiple server back-ends, and a practical benchmarking and testing setup.

The current implementation makes deliberate engineering trade-offs in favor of
simplicity, clarity, and demonstrable functionality. For the intended scope,
those choices are reasonable and produce a coherent end-to-end database system.

\end{document}
